\begin{center}\begin{tikzpicture}[font=\small,>=stealth,line join=round]
\draw[gray!40,rounded corners=3pt] (0,0) rectangle (14.2,8.3);\node at (7.1,7.700000000000001) {例1：奇次项相加减，保留原点对称};
\begin{scope}[shift={(4.1,3.9499999999999997)},x=2.0cm,y=0.3833333333333333cm]\draw[->,gray] (-1.5,0)--(1.5,0) node[right]{$x$};\draw[->,gray] (0,-6)--(0,6) node[above]{$y$};\draw[blue,thick,solid,domain=-1.2:1.2,samples=180] plot(\x,{(\x)^5+3*(\x)^3-2*\x});\end{scope}\node[anchor=west,align=left,text width=5.0cm] at (8.4,4.0) {$H=x^5+3x^3-2x$\\[10pt]三项都随 $x\mapsto-x$ 变号\\[10pt]系数不破坏整体变号\\[10pt]$H(-x)=-H(x)$};
\node[text width=12.6cm,align=center] at (7.1,.65) {实际曲线可有转折；奇偶性讨论对称关系，不要求单调。};
\end{tikzpicture}\end{center}
